% =============================================================================
% CHAPTER 1: Introduction
% =============================================================================
% SPDX-FileCopyrightText: 2026 Frank Langbein <frank@langbein.org>
% SPDX-License-Identifier: LPPL-1.3c
%
% Suggested sections: Motivation and context; Aims and objectives; Contributions;
% Dissertation structure (roadmap). Optional: Scope (boundaries of the work),
% Terminology or key terms. Start with a short chapter intro (no \section) before
% the first section. This file also demonstrates template features; see
% Section~\ref{sec:ch1-reference}. Remove that section for your own report.
% =============================================================================

% -----------------------------------------------------------------------------
% Chapter intro (no section): one short paragraph before the first \section
% -----------------------------------------------------------------------------
This chapter introduces the project's motivation and context for the work, the aims and objectives, the main contributions, and the structure of the remaining chapters. The following sections expand on each of these.

% -----------------------------------------------------------------------------
% Motivation and context (typical first section)
% -----------------------------------------------------------------------------
% Use \section, \subsection, \label{key}. Refer with \ref{key} or \autoref{key}.

\section{Motivation and context}\label{sec:motivation}

% \newthought{This template illustrates} the structure and style of a dissertation using the Qyber dissertation class. As in a real introduction, you would state the problem domain, why it matters, and how your work fits. Acronyms use the long-short style: first use gives the full form (e.g.\@ \gls{uk}); later use \acrshort{api} or \acrlong{api} when you need only short or long.

% As shown in previous work~\citep{example-paper}, we follow standard practice.
% %F Clarify whether this applies to all cases.
% \citet{example-book} provide a longer treatment; multiple citations: \citep{example-paper,example-conf}. You can link to other chapters (e.g.\@ Chapter~\ref{ch2}, Section~\ref{sec:objectives}) and to figures (Fig.~\ref{fig:demo}). The class provides \url{https://example.com} and \href{https://example.com}{link text} for URLs and hyperlinks.

\gls{slam} is the fundamental problem in robotics of an online robot making a map of its surroundings (mapping), whilst simultaneously tracking its own position (localisation). This enables actions such as path planning. This presents a chicken and egg problem, you need the map to calculate position, yet need the position to build the map.

Classical approaches to the \gls{slam} problem represent the environment using only metric terms such as occupancy grids, which captures the raw geometry without any structural meaning. Whilst this is enough for basic navigation, increasingly advanced robotics require a more structured understanding of the scene in order to execute higher-level functions. This is particularly apparent in architecture contexts, where tasks like inventorying or remote visualisation require an understanding of spatial structure.

As acknowledged in \cite{topologic}, current methods for building documentation rely on  inferred geometry and do not explicitly represent spatial connectivity or enclosure. This limits real-world applications such as inventorying, which needs room-level semantic partitioning to associate objects with spaces. Existing solutions such as Topologic \cite{Jabi2025} acknowledges this problem, however it still depends on initial human input for the topology. This motivates an automated system for scanning a real environment to create a topological spatial model, which can then assist in the applications discussed above.

Recent advancements in \gls{slam} research use a graph-based representation of the scene that integrates metric(occupancy) and semantic (classes) estimates, along with hierarchical relationships connecting different aspects of the scene. For example, an object belongs to a room. This approach, known as a \gls{dsg}, enables it to have a broader range of higher-level tasks, such as responding to natural language commands like ‘pickup the red pen’, or architecture-related tasks mentioned above.

Different \gls{dsg} methods use different approaches to how the scene is processed. This means they differ in how the scene is constructed and labelled. This proposes 3 core research questions: Which achieves the best accuracy, which is the most practical for deployment, and which creates the richest semantic understanding? This dissertation addresses these questions through a side-by-side evaluation of 3 state-of-the-art methods, resulting in deploying the best-performing method on real-world data.

% -----------------------------------------------------------------------------
% Aims and objectives (typical second section)
% -----------------------------------------------------------------------------
% Use \citep{} (parenthetical) or \citet{} (narrative), not \cite.

\section{Aims and objectives}\label{sec:objectives}

%State the aim of your project and break it down into objectives. The following is a placeholder list (enumitem package: custom labels).
The main aim of this project was to compare how 3 state-of-the-art methods perform according to different metrics, which would allow a best-performing to be selected to create a \gls{dsg} of the \gls{hccr} laboratory.
The project was broken down into 4 clear objectives:

\begin{enumerate}[label=\textbf{O\arabic*:}, leftmargin=*]
  \item \textbf{Deploy all 3 methods on example data}
    \begin{itemize}
      \item Ensures that the methods are working as intended
    \end{itemize}
  \item \textbf{Deploy methods on custom rosbag and store results}
    \begin{itemize}
      \item Allows a side-by-side performance review between the methods
    \end{itemize}
  \item \textbf{Evaluate results and determine which is the best-performing method}
    \begin{itemize}
      \item Determine the most suitable method for the HCCR lab in terms of metric, semantic and practical performance
    \end{itemize}
  \item \textbf{Use best-performing method on real-world data to generate a \gls{dsg} of the \gls{hccr} laboratory}
    \begin{itemize}
      \item Prove real-world practicality
    \end{itemize}
\end{enumerate}

% When you need mathematics, use amsmath and mathtools. Example:
% \begin{equation}\label{eq:example}
%   E = mc^2, \qquad \int_0^1 f(x)\,\mathrm{d}x = F(1) - F(0).
% \end{equation}
% Refer as Equation~\eqref{eq:example}; inline math: $\alpha + \beta = \gamma$. Paired delimiters (mathtools): \DeclarePairedDelimiter{\abs}{\lvert}{\rvert} then $\abs{x}$, $\abs*{\frac{1}{2}}$. Chemistry (mhchem): \ce{H2O}, \ce{CO2}, \ce{^1H-NMR}. Units (siunitx): \SI{3.0}{m/s}, \num{1234} samples. Symbols: \textdegree, \textcopyright; dingbats (pifont): \ding{51} \ding{43}. Relative size (relsize): \textlarger{larger} \textsmaller{smaller}.

% -----------------------------------------------------------------------------
% Contributions (typical third section)
% -----------------------------------------------------------------------------

\section{Contributions}\label{sec:contributions}

%\epigraph{The best way to predict the future is to invent it.}{Alan Kay}

\begin{itemize}
  \item  \textbf{Integrated RGB-D camera onto existing 3D LiDAR robotic configuration}
    \begin{itemize}
      \item Enabled multimodal data capture
    \end{itemize}
  \item \textbf{Created simulated office dataset using teleoperated robot}
    \begin{itemize}
      \item Created a rosbag recording of scene for reproducable evaluation
    \end{itemize}
  \item \textbf{Deployed 3 state-of-the-art DSG methods (S-Graphs+ \cite{s_graphs_plus}, Hydra \cite{hydra}, Clio \cite{clio}) on the custom rosbag dataset}
    \begin{itemize}
      \item Produced examinable results for evaluation
    \end{itemize}
  \item \textbf{Evaluated methods based on accuracy, practicality, and richness of semantic understanding}
    \begin{itemize}
      \item Enabled a systematic comparison and selection of the best-performing method
    \end{itemize}
  \item \textbf{Deployed best-performing method on real-world data captured in the \gls{hccr} laboratory}
    \begin{itemize}
      \item Generated a \gls{dsg} of the scene and demonstrating the practical application of the method in a real-world context
    \end{itemize}
\end{itemize}
% Summarise the main contributions of the dissertation. You can use a short list and refer to a figure or table. The class provides restyled \texttt{quote} and \texttt{quotation} environments (smaller type, indents):

% \begin{quotation}
%   A short quotation: replace with a real citation or excerpt. Use \verb|\begin{quotation}| \verb|...| \verb|\end{quotation}| for indented, smaller text. For a shorter block use \verb|\begin{quote}| \verb|...| \verb|\end{quote}|.
% \end{quotation}

% \begin{quote}
%   A single-line quote: \verb|\begin{quote}...\end{quote}| gives a shorter block with smaller type (class restyle).
% \end{quote}

% The class also restyles \verb|\begin{verse}...\end{verse}| for short verse or poetry (centred lines, smaller type):
% \begin{verse}
%   Line one.\\
%   Line two.\\
%   Line three.
% \end{verse}

% \begin{itemize}
%   \item First contribution; replace with your own.
%   \item Second contribution.
% \end{itemize}

% Figure~\ref{fig:demo} shows how to include a figure (graphicx, caption; optional short caption for the List of Figures). Float placement: \texttt{[t]} = top of page (default here); \texttt{[htbp]} = here, top, bottom, float page; \texttt{[!t]} = relax limits and prefer top. See the ``Float placement'' comment block in \texttt{main.tex}. \texttt{[H]} (float package) only when you must force ``here'' (can cause uneven spacing). Subfigures (subcaption): see Figure~\ref{fig:sub}.

% \begin{figure}[t]
%   \centering
%   \fbox{\rule{0pt}{4cm}\rule{4cm}{0pt}}
%   \caption[Short for list of figures]{Full caption. Optional argument used in List of Figures.}\label{fig:demo}
% \end{figure}

% \begin{figure}[t]
%   \centering
%   \begin{subfigure}[t]{0.45\linewidth}
%     \centering
%     \fbox{\rule{0pt}{2cm}\rule{2cm}{0pt}}
%     \subcaption{First subfigure.}
%   \end{subfigure}
%   \hfill
%   \begin{subfigure}[t]{0.45\linewidth}
%     \centering
%     \fbox{\rule{0pt}{2cm}\rule{2cm}{0pt}}
%     \subcaption{Second subfigure.}
%   \end{subfigure}
%   \caption{Example with subcaption (a) and (b).}
%   \label{fig:sub}
% \end{figure}

% % Full-width figure: requires class option \texttt{fullwidth}. If you remove that option, comment out this block.
% \begin{fullwidthfigure}[t]
%   \fbox{\rule{0pt}{2.5cm}\rule{\linewidth}{0pt}}
%   \caption[Full-width figure]{Full-width figure example. Use \texttt{\string
%     \begin{fullwidthfigure}[t]}...\texttt{\string
%   \end{fullwidthfigure}} when the class option \texttt{fullwidth} is set.}\label{fig:fullwidth}
% \end{fullwidthfigure}

% % Full-width block (not a float): extends into margin. Requires class option \texttt{fullwidth}; comment out if not set.
% \begin{fullwidth}
%   This paragraph is in \verb|\begin{fullwidth}...\end{fullwidth}|: block content that extends into the left and right margins (e.g.\@ for wide tables or key points).
% \end{fullwidth}

% Tables: use \texttt{booktabs} (\verb|\toprule|, \verb|\midrule|, \verb|\bottomrule|) (Table~\ref{tab:booktabs}), or \texttt{tabularray} with \texttt{tblr} / \texttt{longtable} for flexible or multi-page tables (Table~\ref{tab:demo}, Table~\ref{tab:long}). Same placement options as figures (\texttt{[t]}, \texttt{[!t]}, \texttt{[htbp]}, etc.).

% \begin{table}[t]
%   \centering
%   \caption{Example table (booktabs, Tufte-style).}
%   \label{tab:booktabs}
%   \begin{tabular}{lcc}
%     \toprule
%     Item   & Value A & Value B \\
%     \midrule
%     First  & 1.0     & 2.0    \\
%     Second & 3.0     & 4.0    \\
%     \bottomrule
%   \end{tabular}
% \end{table}

% \begin{table}[t]
%   \centering
%   \caption{Example table (tabularray).}
%   \label{tab:demo}
%   \begin{tblr}{colspec={lcc}, hlines}
%     Item   & Value A & Value B \\
%     First  & 1.0     & 2.0    \\
%     Second & 3.0     & 4.0    \\
%   \end{tblr}
% \end{table}

% \begin{center}
%   \begin{longtable}{lcc}
%     \caption{Minimal longtable demo (not in a float environment).} \label{tab:long} \\
%     \toprule
%     Item   & A & B \\ \midrule
%     \endfirsthead
%     \toprule
%     Item   & A & B \\ \midrule
%     \endhead
%     Row 1  & 1 & 2 \\
%     Row 2  & 3 & 4 \\ \bottomrule
%   \end{longtable}
% \end{center}

% -----------------------------------------------------------------------------
% Dissertation structure / roadmap (typical final section of introduction)
% -----------------------------------------------------------------------------

\section{Dissertation structure}\label{sec:roadmap}

%Outline how the rest of the dissertation is organised so the reader knows what to expect. Reference chapters by label.

The rest of this dissertation is organised as follows.
\begin{itemize}
  \item \textbf{Chapter~\ref{ch2}} (Background): Context, Related Work, Problem Statement.
  \item \textbf{Chapter~\ref{ch3}} (Methodology): Research Design, Specification, Candidate Methods, Comparative Evaluation Framework, Experimental Setup, Implementation.
  \item \textbf{Chapter~\ref{ch4}} (Results and Evaluation): Simulation Output, Results and Discussion, Best-Performing Method Justification, Real-World Results, Limitations.
  \item \textbf{Chapter~\ref{ch5}} (Conclusions and Future Work): Conclusion, Future Work.
  \item \textbf{Chapter~\ref{ch6}} (Reflection on Learning): Critical Narrative, Analysis and Application, Implications for future practice.
  \item \textbf{Appendix~\ref{app:supp}} and \textbf{Appendix~\ref{app:results}}: Supplementary Material and Raw Results.
\end{itemize}

% Algorithms (algorithm, algpseudocode) appear in the List of Algorithms (Algorithm~\ref{alg:demo}). Use \texttt{[t]} or \texttt{[!t]} for top placement like figures/tables. For code snippets use verbatim; for wide tables or figures use \texttt{\string
% \begin{sidewaystable}} or \texttt{\string
%   \begin{sidewaysfigure}} (rotating).\footnote{Footnote example. Captions and footnotes are set at least 11\,pt by the class.} Placeholder prose in other chapters uses \texttt{\string\lipsum[n]}. Example: \lipsum[1]

%     Use \texttt{\string\afterpage\{\string\clearpage\}} (afterpage) when you need to clear after the current page.

%     \begin{algorithm}[t]
%       \caption{Example algorithm (algpseudocode). Alternative: algorithmic for traditional layout.}
%       \label{alg:demo}
%       \begin{algorithmic}[1]
%         \State \textbf{Input:} data $x$
%         \State \textbf{Output:} result $y$
%         \State $y \gets 0$
%         \For{$i = 1$ \textbf{to} $n$}
%         \State $y \gets y + x_i$
%         \EndFor
%         \State \Return $y$
%       \end{algorithmic}
%     \end{algorithm}

%     Short code (verbatim):
% \begin{verbatim}
%   x = 1
%   y = x + 1
% \end{verbatim}

%     % -----------------------------------------------------------------------------
%     % Template reference (for authors: where to find package demos and conventions)
%     % -----------------------------------------------------------------------------

%     \section{Template reference}\label{sec:ch1-reference}

%     \paragraph{Font comparison.}
%     To compare the class font options (\texttt{sans}, \texttt{serif}, \texttt{hacker}), this paragraph uses the default body text, then explicit \textrm{serif}, \textsf{sans}, and \texttt{monospace}, with inline math $E = mc^2$ and a URL \url{https://example.org}. Rebuild with \texttt{font=sans}, \texttt{font=serif}, or \texttt{font=hacker} to see the chosen default; the other families show how headings, code, and captions look in each setup.

%     This introduction doubles as a template: the sections above (Motivation, Aims and objectives, Contributions, Dissertation structure) are the usual structure for a dissertation Chapter~1. Replace their content with your own. The same file also demonstrates:

%     \begin{itemize}
%         \item \textbf{Citations and links:} \verb|\citep|/\verb|\citet| (natbib); link colour DarkBlue, cite/URL colour DarkGreen (class default).
%       \item \textbf{\texttt{\string\newthought}:} Section-opening phrase in small caps with vertical space (e.g.\@ in Motivation). Letterspacing is applied by default.
%         \item \textbf{Quote environments:} \verb|\begin{quotation}| \verb|...| \verb|\end{quotation}|, \verb|\begin{quote}| \verb|...| \verb|\end{quote}|, and \verb|\begin{verse}| \verb|...| \verb|\end{verse}| are restyled (smaller type, indents; verse centred) by the class.
%         \item \textbf{Full-width content:} With class option \texttt{fullwidth}, use \verb|\begin{fullwidth}| \verb|...| \verb|\end{fullwidth}| for block content that extends into the margin, and \texttt{fullwidthfigure}/\texttt{fullwidthtable} for full-width floats (e.g.\@ Figure~\ref{fig:fullwidth}).
%         \item \textbf{Epigraph:} \verb|\epigraph{quote}{source}| for chapter- or section-opening quotes (right-aligned, small type).
%       \item \textbf{Other:} glossaries (long-short acronyms), amsmath/mathtools, textcomp, pifont, siunitx, relsize, mhchem, enumitem, graphicx/caption, float/stfloats, subcaption, tabularray/longtable, algorithm/algpseudocode, verbatim, afterpage, lipsum, rotating. Prefer parenthetical asides in the main text; use footnotes only sparingly (see \texttt{main.tex}).
%     \end{itemize}

%     Conventions for citations, spaces (\verb|~|, \verb|\@|), dashes, and reviewer comments (\verb|%ID |) are in the comment block at the top of \texttt{main.tex}. Class options (parskip/parindent, fullwidth, italic) are documented there too.
